\documentclass[12pt]{article}

% fonts

\usepackage[T1]{fontenc}
\usepackage{newtxtext}
\usepackage{bm}
\usepackage{amsthm} % before newtxmath b/c latex
\usepackage{newtxmath}
% \usepackage[bb=mt,cal=mt]{mathalpha}

\usepackage{amsmath}
\usepackage{amsfonts}
\usepackage{thmtools}
\usepackage{mathtools}

% spacing

\usepackage[parfill]{parskip}
\setlength{\parindent}{0pt}
\setlength{\parskip}{2ex plus 0.2ex minus 0.2ex}

% geometry of the page

\usepackage[
paper  = letterpaper,
left   = 1in,
right  = 1in,
top    = 1in,
bottom = 1in,
]{geometry}

% references

\usepackage{natbib}


\usepackage{tikzit}
\input{pm.tikzstyles}

\begin{document}

\begin{flushleft}
\textbf{%
    An Introduction to Probabilistic Graphical Models (Jordan, 2003) \\
    Chapter 2: Conditional Independence and Factorization
} \\
Satvik Saha (ss7228) \\
\today
\end{flushleft}

\vspace{0.1in}

This chapter introduces a framework for moving back and forth between a
(potentially very complicated) joint probability distribution and its graphical
representation.
One advantage of the latter is that assertions of \emph{conditional
independence} are much easier to read off (almost a triviality in
\emph{undirected} graphical models): the task of carefully factorizing a joint
distribution is translated into examining \emph{graph separations} between
groups of nodes.
This can be done via a \emph{Bayes ball algorithm}: two groups of nodes are
conditionally independent given a third group if `information' cannot travel
along the graph between the former groups, according to a simple set of rules.
In an undirected graphical model, this reduces to checking connectedness
between the former groups after removing the latter!

\vspace{1em}
\ctikzfig{gmm}
\vspace{1em}

For example, in the above (Gaussian mixture) model, we may conclude that each
$z_i \perp z_j \mid \bm{x}, \bm{\mu}, \bm{\lambda}, \pi$ for $i \neq j$, simply
by observing that the paths between $z_i$ and $z_j$ are blocked.
Furthermore, $\pi \perp \{\bm{x}, \bm{\mu}, \bm{\lambda}\} \mid \bm{z}$.

This is often useful when writing full conditionals for Gibbs sampling: the
full conditional of a node depends solely on its \emph{Markov blanket},
consisting of parents, children, and co-parents!
This tells us that the updates for \vspace{-0.5em}
\begin{itemize}
    \itemsep-0.5em
    \item $\pi$ involve only $\bm{z}$,
    \item $z_i$ involve only $x_i, \bm{\mu}, \bm{\lambda}, \pi$,
    \item $\mu_k$ involve only $\bm{x}, \bm{\lambda}, \bm{z}$,
    \item $\lambda_k$ involve only $\bm{x}, \bm{\mu}, \bm{z}$.
\end{itemize}


\nocite{*}
\bibliographystyle{plainnat}

\vfill

\begingroup
\renewcommand{\section}[2]{}%
\textbf{References}
\bibliography{week07}
\endgroup


\end{document}

%%% Local Variables:
%%% mode: latex
%%% TeX-master: t
%%% End:
